\documentclass[10pt]{article}
\usepackage[T1]{fontenc}
\usepackage[utf8]{inputenc}
\usepackage[french]{babel}
\usepackage{amsmath,amssymb,mathrsfs,mathpazo}
\usepackage{geometry}
\geometry{a4paper,margin=2cm}
\usepackage{graphicx,array,tabularx,tabularray,colortbl,multirow,multicol,booktabs,enumitem,pifont,ulem,fancybox,diagbox,cancel}
\usepackage{pstricks,pstricks-add,pst-plot,pst-tree,pst-node,pst-text,pst-func,pst-math,pst-bspline,pst-3dplot,pst-eucl,pst-solides3d}
\let\pscalculate\undefined
\usepackage{pst-bezier}
\usepackage{tikz,pgfplots,tkz-tab}
\usetikzlibrary{arrows,arrows.meta,patterns,calc,positioning,trees,decorations.markings,shapes.geometric,shapes.misc,shapes.arrows,intersections,decorations.pathmorphing,decorations.pathreplacing,backgrounds,shadows,babel}
\pgfplotsset{compat=1.18}
\usepgfplotslibrary{fillbetween}
\usepackage[np]{numprint}
\usepackage{listings,makecell,pas-tableur,twoopt}
\let\barre\undefined
\usepackage[pstricks]{tablvar}
\definecolor{darkgreen}{RGB}{0,100,0}
\providecommand{\red}{\color{red}}
\providecommand{\blue}{\color{blue}}
\providecommand{\green}{\color{green}}
\providecommand{\black}{\color{black}}
\providecommand{\white}{\color{white}}
\providecommand{\gray}{\color{gray}}
\providecommand{\cyan}{\color{cyan}}
\providecommand{\magenta}{\color{magenta}}
\providecommand{\yellow}{\color{yellow}}
\providecommand{\violet}{\color{violet}}
\providecommand{\orange}{\color{orange}}
\providecommand{\euro}{\texteuro}
\providecommand{\EUR}{\texteuro}
\providecommand{\ieme}{\textsuperscript{e}}
\providecommand{\degres}{\ensuremath{^\circ}}
\providecommand{\degree}{\ensuremath{^\circ}}
\providecommand{\vv}[1]{\overrightarrow{#1}}
\providecommand{\Coord}[1]{#1}
\providecommand{\decosix}{}
\providecommand{\starredbullet}{\ensuremath{\bullet}}
\pagestyle{empty}
\setlength{\parindent}{0pt}
\usepackage{adjustbox}\DeclareUnicodeCharacter{2236}{:}\DeclareUnicodeCharacter{2190}{\ensuremath{\leftarrow}}
\providecommand{\R}{\mathbb{R}}
\providecommand{\N}{\mathbb{N}}
\providecommand{\Z}{\mathbb{Z}}
\providecommand{\Q}{\mathbb{Q}}
\providecommand{\D}{\mathbb{D}}
\providecommand{\C}{\mathbb{C}}
\providecommand{\e}{{\mathrm{e}}}
\providecommand{\vect}[1]{\overrightarrow{#1}}
\providecommand{\vectt}[1]{\overrightarrow{\mathrm{#1}}}
\providecommand{\vvt}[1]{\overrightarrow{\mathrm{#1}}}
\providecommand{\np}[2][]{{#2~#1}}
\providecommand{\vv}[1]{{\overrightarrow{#1}}}
\renewcommand{\d}{{\mathrm{d}}}
\providecommand{\degres}{{^\circ}}
\providecommand{\degre}{{^\circ}}
\providecommand{\blue}{\color{blue}}
\providecommand{\red}{\color{red}}
\providecommand{\enskip}{\quad}
\providecommand{\llbracket}{\lbrack\!\lbrack}
\providecommand{\rrbracket}{\rbrack\!\rbrack}
\providecommand{\euro}{{\text{€}}}
\providecommand{\no}{{\text{n°}}}
\providecommand{\interval}[2]{{[#1;#2]}}
\providecommand{\intervFO}[2]{{[#1;#2[}}
\providecommand{\dotfill}{\ldots}

\providecommand{\ds}{\displaystyle}
\providecommand{\textding}[1]{#1}
\providecommand{\Oij}{\(\mathrm{O};\vec{\imath},\vec{\jmath}\)}
\providecommand{\Oijk}{\(\mathrm{O};\vec{\imath},\vec{\jmath},\vec{k}\)}
\providecommand{\Ouv}{\(\mathrm{O};\vec{u},\vec{v}\)}
\providecommand{\intervFF}[2]{\ensuremath{[#1;#2]}}
\renewcommand{\euro}{\ifmmode\text{\texteuro}\else\texteuro\fi}
\DeclareUnicodeCharacter{00B0}{\ifmmode^{\circ}\else\textdegree\fi}
\DeclareUnicodeCharacter{20AC}{\ifmmode\text{\texteuro}\else\texteuro\fi}
% Énoncés : Éducation nationale. Transcription et corrigés : APMEP, Denis Vergès et contributeurs. Découpage et classement : CoopMaths, Sébastien Lozano, Éric Elter et contributeurs. Adaptation de la présentation : Maths Explorer.
% CC BY-SA : https://coopmaths.fr/www/about/
% https://coopmaths.fr/alea/static/bac/2022/tex/bac_2022_05_sujet2_madagascar_1_cor.tex
\begin{document}
\noindent\textit{APMEP / CoopMaths. CC BY-SA. Présentation adaptée par Maths Explorer.}\par\medskip
\medskip

\begin{enumerate}
	\item  L'arbre complété avec les valeurs disponibles :
			\begin{center}
				\pstree[treemode=R,nodesepA=0pt,nodesepB=2.5pt,treesep=1cm,levelsep=2.5cm]{\TR{}}
				{\pstree{\TR{$F$~}\taput{0,48}}
						{\TR{$C$}\taput{0,165}
							\TR{$\overline{C}$}\tbput{0,835}
							}
					\pstree{\TR{$H$~}\tbput{0,52}}
						{\TR{$C$}\taput{0,215}
							\TR{$\overline{C}$}\tbput{0,785}
							}
					}
				\end{center}
	\item Calculons $p(C\cap F)$ : $p(C\cap F)=p_F(C)\times p(F)=0,165\times 0,48=0,0792$
	\item
		\begin{enumerate}
			\item Utilisons la formule des probabilités totales pour calculer $p(C)$ :
			
			$p(C)=p(C\cap F)+p\left(C\cap \overline{F}\right) = p_F(C)\times p(F) + p_{\overline{F}}(C)\times p\left(\overline{F}\right) = 0,165\times 0,48 + 0,215 \times 0,52 = 0,191$.
			\item Si les évènements $F$ et $C$ sont indépendants, alors $p(F\cap C) = p(F)\times  p(C)$.
			
			$p(F\cap C)=0,0792$ \quad et \quad $p(F)\times p(C) = 0,48 \times 0,191=0,09168$.
			
Ces deux résultats sont différents. Les évènements $F$ et $C$ ne sont pas indépendants.
		\end{enumerate}
	\item D'après la formule de Bayes : $p_C(F)=\dfrac{p(F\cap C)}{p(C)}=\dfrac{0,0792}{0,191}\approx 0,4147$.
	
La probabilité de choisir une femme sachant qu'elle est cadre est égale à 0,4147.
	\item 
	\begin{enumerate}
		\item Il s'agit là d'un schéma de Bernoulli : la répétition de 15 expériences aléatoires n'ayant que deux issues, identiques et indépendantes entre elles. Si on note par $X$ la variable aléatoire donnant le nombre de cadres dans l'échantillon, alors $X$ suit donc une loi binomiale de paramètres $n=15$ et $p=0,191$ : $X \sim \mathcal{B}\left(15~;~0,191\right)$
		\item $\begin{gathered}p(X\leqslant 1)=p(X=0)+p(X=1)=\dbinom{15}{0}\times 0,191^0\times (1-0,191)^{15}+\dbinom{15}{1}\times 0,191^1\times (1-0,191)^{15-1}\\ \approx 0,1890\end{gathered}$.
		\item $E(X)=n\times p= 15\times 0,191= 2,865$.
	\end{enumerate}
	\item Soit $n \in \N$.. Cherchons la plus petite valeur de $n$ telle que $p(X\geqslant 1)\geqslant 0,99$.

	$\begin{gathered}p(X\geqslant 1)\geqslant 0,99 \iff  \iff 1-p(X<1)\geqslant 0,99 \iff 1-p(X=0)\geqslant 0,99 \iff -p(X=0)\geqslant -0,01 \\ \iff  p(X=0)\leqslant 0,01 \iff (1-0,191)^n\leqslant 0,01 \iff 0,809^n\leqslant 0,01\end{gathered}$. 
	
La fonction $x\mapsto\ln(x)$ est strictement croissante sur $\R_+^*$ donc
	
	$\begin{gathered}0,809^n\leqslant 0,01 \iff \ln\left(0,809^n \right) \leqslant \ln(0,01) \iff n \times \ln(0,809) \leqslant \ln(0,01) \\ \iff n\geqslant \dfrac{\ln(0,01)}{\ln(0,809)}\end{gathered}$ car $\ln(0,809)<0$. 

À la calculatrice : $\dfrac{\ln(0,01)}{\ln(0,809)}\approx 21,73$ donc $n\geqslant 22$.

Il faudra donc que la taille de l'échantillon choisi soit supérieure ou égale à 22.
\end{enumerate}

\bigskip

\end{document}
